\documentclass[12pt,a4paper]{article}

% ─── Packages ────────────────────────────────────────────────────────────────
\usepackage[top=2.2cm,bottom=2.2cm,left=2.5cm,right=2.5cm]{geometry}
\usepackage[T1]{fontenc}
\usepackage[utf8]{inputenc}
\newcommand{\Rup}{\textbf{Rs.\,}}
\usepackage{lmodern}
\usepackage{microtype}
\usepackage{xcolor}
\usepackage{titlesec}
\usepackage{titling}
\usepackage{graphicx}
\usepackage{booktabs}
\usepackage{longtable}
\usepackage{array}
\usepackage{enumitem}
\usepackage{hyperref}
\usepackage{fancyhdr}
\usepackage{mdframed}
\usepackage{tcolorbox}
\usepackage{multicol}
\usepackage{amsmath}
\usepackage{amssymb}
\usepackage{tabularx}
\usepackage{colortbl}
\usepackage{hhline}
\usepackage{setspace}

% ─── Colour Palette ──────────────────────────────────────────────────────────
\definecolor{primaryGreen}{RGB}{22,101,52}
\definecolor{accentGreen}{RGB}{74,222,128}
\definecolor{deepTeal}{RGB}{15,118,110}
\definecolor{softGold}{RGB}{202,138,4}
\definecolor{lightBg}{RGB}{240,253,244}
\definecolor{cardBg}{RGB}{220,252,231}
\definecolor{borderGreen}{RGB}{134,239,172}
\definecolor{textDark}{RGB}{20,83,45}
\definecolor{sectionBar}{RGB}{16,185,129}
\definecolor{warningOrange}{RGB}{194,65,12}
\definecolor{pilotBlue}{RGB}{37,99,235}

% ─── Typography ──────────────────────────────────────────────────────────────
\titleformat{\section}[block]
  {\large\bfseries\color{primaryGreen}}
  {\thesection.}{0.5em}{}[\vspace{2pt}{\color{sectionBar}\rule{\linewidth}{1.5pt}}]

\titleformat{\subsection}[hang]
  {\normalsize\bfseries\color{deepTeal}}
  {\thesubsection.}{0.5em}{}

\titleformat{\subsubsection}[hang]
  {\small\bfseries\color{textDark}}
  {\thesubsubsection.}{0.5em}{}

\setlength{\parskip}{6pt}
\setlength{\parindent}{0pt}
\onehalfspacing

% ─── Header / Footer ─────────────────────────────────────────────────────────
\pagestyle{fancy}
\fancyhf{}
\fancyhead[L]{\small\color{primaryGreen}\textbf{Kisaan Ki Yash} \textemdash\ 1-km Hyperlocal Agricultural Intelligence Platform}
\fancyhead[R]{\small\color{deepTeal}Smart India Hackathon 2024}
\fancyfoot[C]{\small\color{gray}\thepage}
\renewcommand{\headrulewidth}{0.4pt}
\renewcommand{\headrule}{\hbox to\headwidth{\color{borderGreen}\leaders\hrule height \headrulewidth\hfill}}

% ─── Custom Boxes ────────────────────────────────────────────────────────────
\tcbuselibrary{skins,breakable,most}

\newtcolorbox{kpibox}[2][]{
  colback=lightBg, colframe=primaryGreen, coltitle=white,
  fonttitle=\bfseries\small, title=#2,
  boxrule=1.2pt, arc=5pt, left=8pt, right=8pt, top=6pt, bottom=6pt, #1
}

\newtcolorbox{statbox}{
  colback=cardBg, colframe=sectionBar,
  boxrule=0.8pt, arc=4pt, left=6pt, right=6pt, top=4pt, bottom=4pt
}

\newtcolorbox{warnbox}{
  colback=yellow!10, colframe=softGold,
  boxrule=0.8pt, arc=4pt, left=6pt, right=6pt, top=4pt, bottom=4pt
}

% ─── Hyperref ────────────────────────────────────────────────────────────────
\hypersetup{
  colorlinks=true, linkcolor=primaryGreen, urlcolor=deepTeal,
  pdftitle={Kisaan Ki Yash -- Detailed Overview},
  pdfauthor={Kisaan Ki Yash Team}
}

% ══════════════════════════════════════════════════════════════════════════════
\begin{document}
% ══════════════════════════════════════════════════════════════════════════════

% ─── Title Page ──────────────────────────────────────────────────────────────
\begin{titlepage}
\pagecolor{lightBg}
\centering
\vspace*{2cm}

{\color{primaryGreen}\rule{\linewidth}{2pt}}
\vspace{0.6cm}

{\Huge\bfseries\color{primaryGreen} Kisaan Ki Yash}\\[0.4cm]
{\large\color{deepTeal}\itshape 1-km Hyperlocal Agricultural \& Climate Intelligence Platform}\\[0.3cm]
{\color{sectionBar}\rule{0.5\linewidth}{1pt}}\\[0.8cm]

{\large\bfseries\color{textDark} Detailed Technical \& Business Overview}\\[0.2cm]
{\normalsize\color{gray} Smart India Hackathon 2024 \textemdash\ Problem Statement: Agriculture \& Rural Development}

\vspace{1.5cm}

\begin{tcolorbox}[colback=cardBg,colframe=primaryGreen,arc=8pt,boxrule=1.5pt,width=0.82\linewidth,center]
\centering
\textbf{\color{primaryGreen}Mission Statement}\\[6pt]
\color{textDark}\itshape
``To deliver centimetre-accurate, AI-powered agricultural intelligence directly to every Indian farmer's phone --- translating satellite data, live weather observations, and soil science into plain-language, actionable decisions that protect crops, reduce losses, and increase farm income.''
\end{tcolorbox}

\vspace{1.5cm}

\begin{tabular}{rl}
\textbf{\color{primaryGreen}Platform:}   & Kisaan Ki Yash (``The Glory of the Farmer'') \\
\textbf{\color{primaryGreen}Domain:}     & Agricultural Technology / Climate Intelligence \\
\textbf{\color{primaryGreen}Region:}     & Lucknow District, Uttar Pradesh (Pilot) $\rightarrow$ Pan-India \\
\textbf{\color{primaryGreen}Resolution:} & 1 km $\times$ 1 km hyperlocal grid (EPSG:32644) \\
\textbf{\color{primaryGreen}Models:}     & 10 AI/ML Models (M1--M10) --- 3 Frozen Pilot \\
\textbf{\color{primaryGreen}Platform URL:} & \href{https://kisaankiyash.vercel.app}{kisaankiyash.vercel.app} \\
\end{tabular}

\vfill

{\color{primaryGreen}\rule{\linewidth}{2pt}}
\end{titlepage}

\pagecolor{white}
\newpage
\tableofcontents
\newpage

% ══════════════════════════════════════════════════════════════════════════════
\section{Executive Summary}
% ══════════════════════════════════════════════════════════════════════════════

\textbf{Kisaan Ki Yash} (``The Glory of the Farmer'') is a first-of-its-kind 1-km hyperlocal agricultural intelligence platform built for the Indian farming community. It fuses satellite earth observation data, live AWS (Automatic Weather Station) ground measurements, digital elevation models, and deep-learning downscaling models to produce field-level crop advisory services that no government or commercial platform currently offers at this spatial resolution.

\begin{kpibox}{Platform at a Glance}
\begin{multicols}{3}
\begin{center}
{\Large\bfseries\color{primaryGreen}1 km} \\
\small Spatial resolution grid

\columnbreak

{\Large\bfseries\color{primaryGreen}10} \\
\small AI/ML Intelligence Models

\columnbreak

{\Large\bfseries\color{primaryGreen}\Rup 0--89} \\
\small Monthly access tiers
\end{center}
\end{multicols}
\end{kpibox}

\vspace{4pt}

\begin{kpibox}{Pilot Results (72h, 5 AWS Stations, Lucknow)}
\begin{multicols}{3}
\begin{center}
{\Large\bfseries\color{primaryGreen}39.9\%} \\
\small MAE reduction vs.\ raw NWP forecast

\columnbreak

{\Large\bfseries\color{primaryGreen}0.33\textdegree C} \\
\small Validation MAE (weather downscaling)

\columnbreak

{\Large\bfseries\color{primaryGreen}R\textsuperscript{2} = 0.989} \\
\small Temperature model fit
\end{center}
\end{multicols}
\end{kpibox}

% ──────────────────────────────────────────────────────────────────────────────
\subsection{The Problem}

India has approximately 140 million farm households. Existing advisory services suffer from three fundamental gaps:

\begin{enumerate}[leftmargin=*]
  \item \textbf{Coarse spatial resolution} --- National Weather Service forecasts operate at 25--50 km resolution. A district-wide forecast cannot tell a farmer in Malihabad whether his mango orchard will receive rainfall today or tomorrow.
  \item \textbf{Lack of personalisation} --- Generic advisories issued over All India Radio or Kisan Call Centre are not calibrated to soil type, crop stage, micro-topography, or individual farm location.
  \item \textbf{Reactive, not predictive} --- Farmers currently respond after damage occurs (pest outbreak, waterlogging, frost). There is no proactive, location-specific system that tells them to act \textit{before} a risk materialises.
\end{enumerate}

% ──────────────────────────────────────────────────────────────────────────────
\subsection{Our Solution}

Kisaan Ki Yash solves all three gaps simultaneously:

\begin{itemize}[leftmargin=*]
  \item \textbf{1-km hyperlocal grid} --- we downscale 25-km NWP forecasts to 1-km using Random Forest + terrain features verified against 5 AWS ground stations.
  \item \textbf{Panchayat-level personalisation} --- each prediction is tagged to the farmer's specific panchayat, soil profile, and crop calendar.
  \item \textbf{Proactive, multi-day intelligence} --- alerts are pushed 48--72 hours in advance so farmers can act preventively, not reactively.
\end{itemize}

% ══════════════════════════════════════════════════════════════════════════════
\section{Technical Architecture}
% ══════════════════════════════════════════════════════════════════════════════

\subsection{System Overview}

The platform follows a sequential, dependency-aware intelligence pipeline composed of 10 specialised AI/ML models (M1 through M10), a FastAPI REST backend, and a React/TypeScript frontend.

\begin{center}
\small
\begin{tabular}{|c|p{5.2cm}|p{2.5cm}|p{3.5cm}|}
\hline
\rowcolor{primaryGreen}
\textbf{\color{white}ID} & \textbf{\color{white}Model Name} & \textbf{\color{white}Status} & \textbf{\color{white}Output} \\
\hline
M1 & Hyperlocal Weather Downscaling & \textcolor{primaryGreen}{\textbf{FROZEN PILOT}} & 1-km Temperature, RH, Pressure \\
\hline
M2 & High-Res Temperature Refinement & \textcolor{primaryGreen}{\textbf{FROZEN PILOT}} & Micro-climate Temperature \\
\hline
M3 & Precipitation Downscaling (Hurdle) & \textcolor{primaryGreen}{\textbf{FROZEN PILOT}} & Rain Probability, Intensity \\
\hline
M4 & Soil Moisture Intelligence & \textcolor{pilotBlue}{In Development} & Root-zone Moisture \\
\hline
M5 & Crop Phenology Intelligence & \textcolor{pilotBlue}{In Development} & Growth Stage, GDD Tracking \\
\hline
M6 & ET \& Irrigation Demand & \textcolor{pilotBlue}{In Development} & Daily Irrigation Schedule \\
\hline
M7 & Crop Yield Forecasting & \textcolor{pilotBlue}{In Development} & End-of-season Yield Estimate \\
\hline
M8 & Flood \& Waterlogging Risk & \textcolor{pilotBlue}{In Development} & Inundation Risk Map \\
\hline
M9 & Extreme Weather \& Hazard & \textcolor{pilotBlue}{In Development} & Heatwave / Frost / Dry Spell Alerts \\
\hline
M10 & Agricultural Decision Engine & \textcolor{pilotBlue}{In Development} & Advisory + Expected-Loss Engine \\
\hline
\end{tabular}
\end{center}

\subsection{Data Sources}

\begin{itemize}[leftmargin=*]
  \item \textbf{NWP (Numerical Weather Prediction):} NCMRWF/IMD 0.25\textdegree  gridded forecasts
  \item \textbf{AWS Observations:} 5 calibrated ground stations in Lucknow (Amausi, BKT, Chinhat, Mohanlalganj, Malihabad)
  \item \textbf{Digital Elevation Model:} SRTM 30m $\rightarrow$ aggregated to 1-km grid
  \item \textbf{Satellite:} Sentinel-1 SAR (M4), Sentinel-2 multispectral (M5), MODIS LST
  \item \textbf{Reanalysis:} ERA5-Land hourly (historical feature context)
  \item \textbf{Soil:} ISRIC SoilGrids 250m (clay, sand, bulk density, SOC)
  \item \textbf{Administrative:} Panchayat, Block, District boundary polygons
\end{itemize}

\subsection{Spatial Framework}

All models operate on a standardised 1 km $\times$ 1 km metric grid in \textbf{EPSG:32644} (UTM Zone 44N), with WGS84 (EPSG:4326) used at API input/output boundaries. This eliminates distortion in area calculations and ensures consistent spatial joins across terrain, satellite, and administrative layers.

\subsection{Backend Architecture}

\begin{itemize}[leftmargin=*]
  \item \textbf{REST API:} FastAPI (Python) with full Pydantic typed schemas
  \item \textbf{Model Governance:} ONE\_ACTIVE\_TRAINING\_JOB invariant --- exactly one ML training job at a time; frozen models cryptographically verified by SHA-256 on every startup
  \item \textbf{Uncertainty Quantification:} Conformal prediction intervals (split-conformal residual quantile) for M1/M2; honest NULL bounds returned for M3 (limited calibration)
  \item \textbf{OOD Gating:} Any spatial query outside the validated Lucknow pilot domain returns ABSTAIN status --- no fabricated confidence
  \item \textbf{Provenance Tracking:} Every inference stores dataset version, artifact hash, feature list, baseline reference, and agent ID in an immutable audit log
\end{itemize}

\subsection{Frontend}

\begin{itemize}[leftmargin=*]
  \item React 18 + TypeScript + Vite
  \item Cinematic hero with live-animated climate orb
  \item M1 weather + M3 precipitation dual-mode forecast card
  \item Pipeline stepper showing KisanIntelligenceCore model flow
  \item Live at: \href{https://kisaankiyash.vercel.app}{kisaankiyash.vercel.app}
\end{itemize}

% ══════════════════════════════════════════════════════════════════════════════
\section{Features --- Implemented \& Planned}
% ══════════════════════════════════════════════════════════════════════════════

\subsection{Implemented (Pilot Phase)}

\begin{kpibox}{Model 1 --- Hyperlocal Weather Downscaling}
\begin{itemize}[leftmargin=*,noitemsep]
  \item Random Forest trained on 3 AWS stations; tested on 1 fully held-out station
  \item 39.9\% improvement in MAE over raw NWP (0.408\textdegree C vs.\ 0.680\textdegree C)
  \item Conformal prediction intervals at 90\% nominal coverage
  \item LOSO (Leave-One-Station-Out) spatial holdout --- zero data leakage
  \item Features: coarse T2m, RH, pressure, elevation, slope, aspect, cropland fraction, solar geometry (sin/cos hour, sin/cos DOY)
\end{itemize}
\end{kpibox}

\begin{kpibox}{Model 2 --- Temperature Micro-Climate Refinement}
\begin{itemize}[leftmargin=*,noitemsep]
  \item Refines M1 output with canopy/solar geometry correction
  \item Locked-test MAE 0.411\textdegree C; consistent with M1 (expected at this pilot scale)
  \item Terrain-adjusted lapse rate applied per 1-km cell
\end{itemize}
\end{kpibox}

\begin{kpibox}{Model 3 --- Precipitation Downscaling (Two-Stage Hurdle)}
\begin{itemize}[leftmargin=*,noitemsep]
  \item Stage 1: Logistic Regression --- rain occurrence probability
  \item Stage 2: Ridge regression on log1p --- conditional intensity (mm/h)
  \item Gated hurdle output: expected rainfall = P(rain) $\times$ conditional intensity
  \item Honest uncertainty: bounds withheld when calibration data is insufficient
\end{itemize}
\end{kpibox}

\subsection{Features Being Integrated Before Final Submission}

\subsubsection{M4 --- Soil Moisture Intelligence}
Fusion of Sentinel-1 C-band SAR backscatter with SoilGrids 250m properties to estimate surface and root-zone soil moisture at 1-km resolution. Enables direct irrigation scheduling without in-field sensors.

\subsubsection{M5 --- Crop Phenology \& Growth Stage Tracking}
NDVI and EVI time series (Sentinel-2) classified into growth stages (emergence $\rightarrow$ tillering $\rightarrow$ heading $\rightarrow$ maturity) and mapped to growing degree day (GDD) accumulation calendars. Specific to paddy, wheat, mustard, and mango orchards.

\subsubsection{M6 --- ET and Irrigation Demand Engine}
Full FAO-56 Penman-Monteith evapotranspiration computation using M1/M2 weather inputs. Outputs daily irrigation requirement in mm, integrated with soil moisture deficit (M4) and crop coefficient (Kc) from M5 growth stage.

\subsubsection{M7 --- Crop Yield Forecasting}
End-of-season yield estimate conditioned on cumulative weather anomalies (departure from 30-year normal), vegetation stress index, and irrigation history. District-block calibrated against official CCE (Crop Cutting Experiment) yield data.

\subsubsection{M8 --- Flood \& Waterlogging Risk Intelligence}
Topographic Wetness Index (TWI) and DEM depression analysis to identify waterlogging-prone 1-km cells. M3 precipitation forecast drives a simple water balance to issue 24--72 hour inundation risk alerts.

\subsubsection{M9 --- Extreme Weather \& Agronomic Hazard Detection}
Climatology percentile thresholds (30-year ERA5 baseline) to detect:
\begin{itemize}[leftmargin=*, noitemsep]
  \item Heatwave events (>95th percentile temperature for 3+ consecutive days)
  \item Cold wave / frost risk (temperature $<$ 4\textdegree C during rabi season)
  \item Prolonged dry spell (zero rainfall $>$ 7 days during kharif)
  \item Cyclone/severe storm proximity alerts
\end{itemize}

\subsubsection{M10 --- Agricultural Decision Intelligence Engine}
Rule engine + expected-loss optimisation that aggregates M1--M9 outputs into a single \textbf{Master Advisory} per panchayat. Actions include:
\begin{itemize}[leftmargin=*, noitemsep]
  \item Irrigation: RECOMMENDED / CONDITIONAL / PROHIBITED
  \item Pesticide/Fertiliser Spraying: timing window with rainfall-safe check
  \item Sowing: go/no-go based on soil temperature and moisture
  \item Harvest: risk of rainfall during drying period
  \item Hazard Mitigation: flood berm preparation, frost cover advisory
\end{itemize}

\subsubsection{Weather-Triggered Fertiliser Alert System}
\textit{(New Feature --- Pre-Submission Integration)}

A specialised sub-module of M10 that monitors 7-day precipitation forecasts and triggers SMS/in-app push notifications:

\begin{statbox}
\textbf{How it works:}\\
``Rain is expected in your region (Malihabad, Lucknow) in 5--7 days. This is the optimal window to apply your basal dose of DAP fertiliser. Apply before \textbf{October 3rd} for best root uptake. Avoid application within 48 hours of expected rainfall.''
\end{statbox}

This single feature can prevent an estimated \Rup 2,000--8,000/acre loss per season due to mis-timed fertiliser application (nutrient runoff during heavy rain, or poor uptake in dry soil).

\subsubsection{Multi-Language Vernacular Advisory}
All advisories generated in Hindi (\textit{Devanagari}), English, and transliterated Hinglish for voice-forward UX compatibility with feature phones.

\subsubsection{IVR Voice Alert Integration}
For non-smartphone farmers, outbound IVR (Interactive Voice Response) calls in local language deliver the same advisory that smart-phone users receive via push notification.

\subsubsection{PMFBY / e-Nam Integration Layer}
Verified yield forecasts (M7) and weather event reports (M9) auto-generate farmer-ready documentation for PMFBY (Pradhan Mantri Fasal Bima Yojana) crop insurance claims and price discovery on e-NAM.

% ══════════════════════════════════════════════════════════════════════════════
\section{Business Model}
% ══════════════════════════════════════════════════════════════════════════════

\subsection{Revenue Philosophy}

Kisaan Ki Yash operates on a \textbf{freemium + micro-subscription} model. The base weather and rain alert is \textbf{free forever} --- ensuring digital inclusion. Premium advisory tiers are priced at levels that any Indian farmer can afford, lower than the cost of a single failed irrigation event.

\subsection{Pricing Tiers}

\begin{center}
\small
\renewcommand{\arraystretch}{1.5}
\begin{tabular}{|p{3cm}|p{2.5cm}|p{7.5cm}|}
\hline
\rowcolor{primaryGreen}
\textbf{\color{white}Tier} & \textbf{\color{white}Price} & \textbf{\color{white}What's Included} \\
\hline
\textbf{Kisan Free} & \textbf{\Rup 0 / month} & Hyperlocal temperature, daily rain alert, 3-day basic forecast \\
\hline
\textbf{Kisan Saathi} & \textbf{\Rup 49 / month} & All Free + Soil moisture, irrigation schedule, fertiliser timing alerts (SMS + app), 7-day forecast \\
\hline
\textbf{Kisan Pro} & \textbf{\Rup 89 / month} & All Saathi + Crop yield forecast, pest/disease risk, extreme weather hazard alerts, PMFBY report export, Hindi voice advisory calls \\
\hline
\textbf{Agri-Business / FPO} & \textbf{\Rup 499 / month} & Multi-farm dashboard, API access, panchayat aggregated risk reports, export to Excel/PDF \\
\hline
\textbf{Institutional} & \textbf{Custom} & State government, ATMA, KVK integration, white-label, bulk SMS gateway \\
\hline
\end{tabular}
\end{center}

\subsection{Revenue Streams}

\begin{enumerate}[leftmargin=*]
  \item \textbf{Direct Subscription (B2C):} \Rup 49--89/month from individual farmers via UPI / Aadhaar-linked payment
  \item \textbf{FPO/Co-operative Licensing (B2B):} Farmer Producer Organisations license the platform for member farms at group discount
  \item \textbf{Government API Licensing (B2G):} State Agriculture Departments (ATMA, KVK) integrate our API into their existing digital infrastructure
  \item \textbf{Agri-Input Company Partnerships:} Fertiliser, pesticide, and seed companies pay for precision placement of product advisories (not spam --- weather-triggered and agronomically justified)
  \item \textbf{Crop Insurance Data:} Verified micro-weather event reports sold as evidence for PMFBY claim validation (replaces subjective field inspection)
  \item \textbf{Carbon Credit Verification:} Precision irrigation scheduling reduces water usage and GHG emissions --- certifiable for voluntary carbon markets
\end{enumerate}

\subsection{Unit Economics}

\begin{statbox}
\textbf{Break-Even Analysis (Conservative):}\\
At \Rup 59/month average subscription $\times$ 17,000 paying farmers = \textbf{\Rup 10 lakh/month} gross revenue \\
Target market: Lucknow district alone has 2.5 lakh farm households \\
\textbf{Payback period at 1\% penetration:} $<$ 6 months from pilot validation
\end{statbox}

\subsection{Market Size}

\begin{itemize}[leftmargin=*]
  \item \textbf{TAM:} 140 million Indian farm households --- \Rup 82,000 crore addressable agri-tech market
  \item \textbf{SAM:} 15 million digitally active farmers in Uttar Pradesh + Punjab + Maharashtra
  \item \textbf{SOM (Year 1):} 50,000 farmers in Lucknow + Hardoi + Unnao districts
\end{itemize}

\subsection{Go-to-Market Strategy}

\begin{enumerate}[leftmargin=*, noitemsep]
  \item \textbf{SIH Win $\rightarrow$ Government Validation:} Use hackathon outcome as institutional credibility stamp
  \item \textbf{KVK Partnerships:} Embed into existing Krishi Vigyan Kendra extension network (free adoption for farmers)
  \item \textbf{WhatsApp Bot:} Zero-app-install distribution --- farmers receive advisory via WhatsApp in Hindi
  \item \textbf{Micro-agent Network:} Commission-based local agents (like Business Correspondents for banking) onboard farmers in each village
  \item \textbf{FPO Anchor Deals:} Negotiate block deals with 5--10 large Farmer Producer Organisations
\end{enumerate}

% ══════════════════════════════════════════════════════════════════════════════
\section{Scientific Validation \& Pilot Integrity}
% ══════════════════════════════════════════════════════════════════════════════

All models follow a strict scientific governance protocol:

\begin{itemize}[leftmargin=*]
  \item \textbf{LOSO Holdout:} One entire weather station (AWS\_LKO\_05, Malihabad) never seen during training or hyperparameter tuning --- used only for final locked-test evaluation
  \item \textbf{No Data Leakage:} Future timestamps strictly excluded from feature windows; temporal CV never bleeds forward
  \item \textbf{No Fabricated Metrics:} Conformal bounds reported as NULL when calibration data is insufficient --- system abstains rather than guesses
  \item \textbf{SHA-256 Artifact Locking:} Model weights are cryptographically locked after validation; any tampering detected on boot
  \item \textbf{Limitation Disclosure:} Current pilot covers 72 hours $\times$ 5 stations --- explicitly labelled PILOT ONLY in all interfaces
\end{itemize}

% ══════════════════════════════════════════════════════════════════════════════
\section{Social Impact}
% ══════════════════════════════════════════════════════════════════════════════

\begin{kpibox}{Projected Impact (Year 1 --- 50,000 Farmers)}
\begin{multicols}{2}
\textbf{Economic:}
\begin{itemize}[leftmargin=*,noitemsep]
  \item \Rup 4,000--12,000/farmer/year savings from optimised inputs
  \item 15--25\% reduction in irrigation water usage
  \item 20\% reduction in fertiliser runoff loss
  \item 8--15\% crop loss prevention from early hazard alerts
\end{itemize}

\columnbreak

\textbf{Environmental:}
\begin{itemize}[leftmargin=*,noitemsep]
  \item 30\% reduction in over-irrigation (water table protection)
  \item Precision fertiliser timing $\rightarrow$ 20\% lower N\textsubscript{2}O emissions
  \item Enables verifiable carbon credit generation per farm
  \item Satellite-verified crop insurance $\rightarrow$ prevents false claims
\end{itemize}
\end{multicols}
\end{kpibox}

\begin{kpibox}{SDG Alignment}
\begin{multicols}{2}
\begin{itemize}[leftmargin=*,noitemsep]
  \item \textbf{SDG 1} --- No Poverty: farm income protection
  \item \textbf{SDG 2} --- Zero Hunger: yield stabilisation
  \item \textbf{SDG 6} --- Clean Water: precision irrigation
  \item \textbf{SDG 13} --- Climate Action: emission reduction
  \item \textbf{SDG 15} --- Land: sustainable soil management
  \item \textbf{SDG 9} --- Industry/Innovation: rural AI infrastructure
\end{itemize}
\end{multicols}
\end{kpibox}

% ══════════════════════════════════════════════════════════════════════════════
\section{Competitive Landscape}
% ══════════════════════════════════════════════════════════════════════════════

\begin{center}
\small
\renewcommand{\arraystretch}{1.4}
\begin{tabular}{|p{3.5cm}|p{2cm}|p{2cm}|p{2cm}|p{2.5cm}|}
\hline
\rowcolor{primaryGreen}
\textbf{\color{white}Platform} & \textbf{\color{white}Resolution} & \textbf{\color{white}Ground Truth} & \textbf{\color{white}Offline Access} & \textbf{\color{white}India Price} \\
\hline
Kisaan Ki Yash & \textbf{1 km} & \textbf{5 AWS stations} & \textbf{Yes (SMS)} & \textbf{Free--\Rup 89} \\
\hline
Skymet Agriculture & 25 km & Limited & No & \Rup 300+ \\
\hline
IBM PAIRS (Watson) & 1 km & Enterprise & No & Enterprise \\
\hline
Fasal & Field sensors & In-field IoT & Partial & \Rup 12,000 device \\
\hline
Cropin & Satellite & No ground & No & Enterprise \\
\hline
IMD Mausam App & District & 0 ground & No & Free (no farm AI) \\
\hline
\end{tabular}
\end{center}

\textbf{Key Differentiator:} Kisaan Ki Yash is the only platform combining 1-km physically downscaled NWP with AWS ground truth calibration, soil + vegetation satellite fusion, conformal uncertainty quantification, \textbf{and} vernacular push alerts at a micro-subscription price accessible to small \& marginal farmers.

% ══════════════════════════════════════════════════════════════════════════════
\section{Roadmap}
% ══════════════════════════════════════════════════════════════════════════════

\begin{center}
\small
\renewcommand{\arraystretch}{1.5}
\begin{tabular}{|p{2.5cm}|p{4cm}|p{7cm}|}
\hline
\rowcolor{primaryGreen}
\textbf{\color{white}Timeline} & \textbf{\color{white}Milestone} & \textbf{\color{white}Deliverables} \\
\hline
\textbf{Q4 2024} & SIH Submission & M1--M3 frozen pilot, backend API, live website, WhatsApp bot MVP \\
\hline
\textbf{Q1 2025} & Pilot Expansion & ERA5-Land historical integration, multi-season data, M4 soil moisture \\
\hline
\textbf{Q2 2025} & Beta Launch & M5/M6 irrigation engine, 500 farmer beta users, KVK partnership \\
\hline
\textbf{Q3 2025} & Commercial Launch & All 10 models, UPI billing, FPO onboarding, IVR voice alerts \\
\hline
\textbf{Q1 2026} & Pan-India Scale & PMFBY integration, e-NAM data layer, 5 state rollout \\
\hline
\end{tabular}
\end{center}

% ══════════════════════════════════════════════════════════════════════════════
\section{Team \& Technology Stack}
% ══════════════════════════════════════════════════════════════════════════════

\subsection{Technology Stack}

\begin{multicols}{2}
\textbf{ML/Data Science:}
\begin{itemize}[leftmargin=*,noitemsep]
  \item Python 3.14, scikit-learn, XGBoost
  \item numpy, pandas, joblib
  \item Conformal prediction (split-residual)
  \item GDAL, rasterio (geospatial)
\end{itemize}

\columnbreak

\textbf{Backend:}
\begin{itemize}[leftmargin=*,noitemsep]
  \item FastAPI, Pydantic v2, Uvicorn
  \item REST-first, typed schemas
  \item SHA-256 model registry
  \item Audit log + telemetry metrics
\end{itemize}
\end{multicols}

\begin{multicols}{2}
\textbf{Frontend:}
\begin{itemize}[leftmargin=*,noitemsep]
  \item React 18, TypeScript, Vite
  \item TailwindCSS + custom glassmorphism
  \item Three.js climate orb animation
  \item Deployed: Vercel
\end{itemize}

\columnbreak

\textbf{Data Infrastructure:}
\begin{itemize}[leftmargin=*,noitemsep]
  \item EPSG:32644 UTM metric grid
  \item Sentinel-1/2, ERA5-Land, SRTM
  \item IMD AWS observation network
  \item SoilGrids 250m (ISRIC)
\end{itemize}
\end{multicols}

% ══════════════════════════════════════════════════════════════════════════════
\section*{Conclusion}
\addcontentsline{toc}{section}{Conclusion}
% ══════════════════════════════════════════════════════════════════════════════

Kisaan Ki Yash represents a convergence of modern machine learning, satellite earth observation, and behavioural economics to solve one of India's oldest and most persistent challenges: the disconnect between atmospheric science and the farmer who stands in a field, making decisions with incomplete information.

By grounding every prediction in real AWS observations, maintaining strict scientific honesty about what the models can and cannot claim, and delivering intelligence at a price point that any farmer can afford, this platform has the potential to measurably improve the lives of millions of farm households --- starting from a single panchayat in Lucknow and scaling to every corner of India.

\vspace{0.5cm}

\begin{center}
{\large\bfseries\color{primaryGreen} Kisaan Ki Yash --- When science reaches the field, that is the farmer's glory.}\\
\textit{When science reaches the field --- that is the farmer's glory.}
\end{center}

\end{document}
